% Lösungen Einheit 2 von 3 – Kreisfläche (zusammenbau v0.1)
\einheitenkopf{Einheit 2 von 3 · Kreisfläche}

%% TODO Lösung der Erklärzeile 15g) fehlt in der Bank
\erg{15}{a) $A = \pi \cdot r^2$ \quad b) $A = \pi \cdot 13^2 \approx 530{,}9$ cm² \quad c) $A = \pi \cdot 9^2 \approx 254{,}5$ m² \quad d) $A = \pi \cdot 7^2 \approx 153{,}9$ cm² \quad e) $A = \pi \cdot 16^2 \approx 804{,}2$ mm² \quad f) $A = \pi \cdot 15^2 \approx 706{,}9$ cm² \quad g) \ldots \quad h) $r = 8$ cm, $A = \pi \cdot 8^2 \approx 201{,}1$ cm² \quad i) $A = \pi \cdot 2{,}6^2 \approx 21{,}24$ m² \quad j) $A = \pi \cdot 6^2 : 2 \approx 56{,}5$ m² \quad k) $d = 7$ cm, $u \approx 22{,}0$ cm, $A \approx 38{,}5$ cm² \quad l) $r = \sqrt{154 : \pi} \approx 7{,}0$ cm \quad m) $r = 1{,}18$ m, $A = \pi \cdot 1{,}18^2 \approx 4{,}37$ m²}
\erg{16}{a) $\frac{1}{4} \cdot \pi \cdot 7^2 \approx 38{,}5$ cm²}
\erg{17}{a) $r = 14$ cm, $A = \pi \cdot 14^2 \approx 616$ cm²}
\erg{18}{a) Umfang: $u = 2 \cdot \pi \cdot 3 \approx 18{,}8$ m}
\erg{19}{a) Max hat den Durchmesser eingesetzt statt des Radius. Richtig: $r = 6$ cm, $A = \pi \cdot 6^2 \approx 113{,}1$ cm² \quad b) Die Bögen der Stücke bilden oben und unten den Rand; die Hälfte der Bögen liegt oben, die andere Hälfte unten – jede lange Seite ist also der halbe Umfang. \quad c) groß: $\pi \cdot 18^2 \approx 1\,017{,}9$ cm², $1\,017{,}9 : 12 \approx 84{,}8$ cm² je €; klein: $2 \cdot \pi \cdot 13^2 \approx 1\,061{,}9$ cm², $1\,061{,}9 : 15 \approx 70{,}8$ cm² je €; die große Pizza lohnt sich mehr. \quad d) Halbkreis mit $r = 3$}
